\section{Base de Datos: Preparación para Entornos de Producción}
\label{sec:s5_database}

La preparación para producción combina la siembra de datos maestros y una política de respaldos
automatizada. La Figura~\ref{fig:s5_backup} resume el ciclo de respaldo y recuperación sobre
la base de datos PostgreSQL~15.10 autoalojada en el servidor TIS.\par

\begin{figure}[!ht]
    \centering
    \begin{tikzpicture}[
        node distance=0.6cm,
        s/.style={draw=colorPrimario, fill=headerTabla, thick, rounded corners=2pt,
            text width=2.4cm, align=center, minimum height=1.0cm, font=\sffamily\scriptsize},
        rel/.style={-{Latex[length=2mm]}, colorPrimario, font=\sffamily\tiny}
    ]
        \node[s] (a) {1. Backup\\programado};
        \node[s, right=of a] (b) {2. Almacén\\seguro};
        \node[s, right=of b] (c) {3. Verifica\\integridad};
        \node[s, right=of c] (d) {4. Restaura\\(contingencia)};
        \draw[rel] (a)--(b); \draw[rel] (b)--(c); \draw[rel] (c)--(d);
    \end{tikzpicture}
    \caption{Ciclo de respaldo y recuperación de datos (Sprint 5).}
    \label{fig:s5_backup}
\end{figure}

\subsection{Inyección de Semillas de Datos Maestros (Seeders)}
\label{ssec:s5_db_seeders}
Se inyectaron las semillas de datos maestros requeridas para el arranque operativo: catálogos
base, configuraciones por defecto y datos de referencia, garantizando que el sistema sea
funcional desde el primer despliegue.

\subsection{Políticas de Respaldos y Recuperación}
\label{ssec:s5_db_backups}
Se implementaron políticas automatizadas de respaldos (\textit{backups}) sobre la base de datos
PostgreSQL~15.10 (volcados con \comando{pg\_dump}) y protocolos de recuperación ante
contingencias, definiendo la frecuencia de copia y el procedimiento de restauración.

\begin{NotaTecnica}
Las migraciones se aplican sobre PostgreSQL~15.10 de forma controlada a través del cliente web
phpPgAdmin (\url{http://bytebusters.tis.cs.umss.edu.bo/phppgadmin}); cada script se revisa antes
de ejecutarlo en un entorno compartido para preservar la integridad del schema \comando{core}.
\end{NotaTecnica}

\subsection{Catálogo de Semillas y Política de Respaldos}
\label{ssec:s5_db_catalogo}
La tabla~\ref{tab:s5_seeds} resume los datos maestros sembrados y la política de respaldos
aplicada para el arranque operativo.

\begin{table}[!ht]
    \centering
    \renewcommand{\arraystretch}{1.3}
    \fontsize{9}{10.8}\selectfont\sffamily
    \begin{tabularx}{\textwidth}{|l|X|}
        \hline
        \rowcolor{colorPrimario}
        \textcolor{white}{\textbf{Elemento}} & \textcolor{white}{\textbf{Contenido / criterio}} \\ \hline
        \comando{global\_skill\_tags} & Catálogo base de \textit{skills} con categorías. \\ \hline
        \rowcolor{grisTecnico}
        Cuentas administrativas & Semilla de cuentas \comando{adm.*@ethoshub.dev}. \\ \hline
        Configuración por defecto & Valores iniciales de \comando{portfolio\_settings}. \\ \hline
        \rowcolor{grisTecnico}
        Respaldo programado & Copia periódica con \comando{pg\_dump} del servidor TIS. \\ \hline
        Verificación de integridad & Comprobación del respaldo antes de archivarlo. \\ \hline
        \rowcolor{grisTecnico}
        Recuperación & Procedimiento de restauración ante contingencia. \\ \hline
    \end{tabularx}
    \caption{Semillas maestras y política de respaldos (Sprint 5).}
    \label{tab:s5_seeds}
\end{table}

El ciclo completo de respaldo y recuperación se modela en la Figura~\ref{fig:s5_backup},
garantizando la continuidad operativa ante incidentes en producción.

\clearpage
